\section{Version Control}

Git is used for version control of documentation, infrastructure
configuration, scripts, and application code.

The primary repository is:

\begin{center}
\texttt{homelab-infrastructure}
\end{center}

\section{Main Branch}

The primary branch is:

\begin{center}
\texttt{main}
\end{center}

Direct development on the main branch should be avoided.

\section{Branching Strategy}

Changes are developed using dedicated branches.

Examples include:

\begin{verbatim}
feature/host-only-network
feature/docker-environment
docs/project-architecture
fix/network-configuration
\end{verbatim}

Completed work is integrated into the \texttt{main} branch using Pull Requests.

\newpage
\section{Commit Convention}

Commit messages use semantic prefixes such as:

\begin{verbatim}
feat: Introduces a new feature, service, infrastructure component, or significant 
      new capability.
fix: Corrects an error, misconfiguration, or unexpected behaviour in an existing 
     component.
docs: Introduces changes that affect documentation without changing the behaviour 
      of the infrastructure or application.
refactor: Changes the internal structure of code or configuration without 
      introducing a new feature or fixing a specific defect.
test: Adds or modifies tests, validation procedures, or automated verification 
      mechanisms.
chore: Represents maintenance work that does not directly modify application or 
       infrastructure functionality.
\end{verbatim}

Commits should represent logically coherent changes rather than large,
unrelated batches of modifications.